\documentclass[11pt, a4paper]{article}
\usepackage[utf8]{inputenc}
\usepackage{amsmath, amssymb}
\usepackage{graphicx}
\usepackage{hyperref}
\usepackage{xcolor}
\usepackage{listings}
\lstset{
    language=Python,
    basicstyle=\ttfamily\small,
    keywordstyle=\color{blue}\bfseries,
    commentstyle=\color{gray}\itshape,
    stringstyle=\color{orange},
    numberstyle=\tiny\color{gray},
    numbers=left,
    stepnumber=1,
    numbersep=8pt,
    frame=single,
    breaklines=true,
    breakatwhitespace=false,
    showstringspaces=false,
    tabsize=4,
    captionpos=b
}
\usepackage{geometry}
\geometry{margin=1in}

\title{EE 325 Digital Signal Processing \\ Assignment 6: Realization of Discrete-Time Systems}
\author{PERERA J.D.T. \\ E/21/291}
\date{\today}

\begin{document}
\maketitle

\section{1. Direct Form I Realization}
\textbf{Difference Equation:}
$$y[n] - 1.1y[n-1] + 0.3y[n-2] = x[n] + 0.5x[n-1]$$
$$y[n] = 1.1y[n-1] - 0.3y[n-2] + x[n] + 0.5x[n-1]$$

\textbf{Block Diagram (Direct Form I):}
Direct Form I implements the feedforward part ($x[n]$ terms) and feedback part ($y[n]$ terms) separately, connected in cascade.

Specifically for $y[n] = 1.1y[n-1] - 0.3y[n-2] + ...$:
\begin{itemize}
    \item Tap from $y[n-1]$ is multiplied by $1.1$.
    \item Tap from $y[n-2]$ is multiplied by $-0.3$.
\end{itemize}

\textbf{Delay Elements:} 3 total ($1$ for $x$, $2$ for $y$). Note: This is non-canonic.

\section{2. Direct Form II Realization}
\textbf{(a) Direct Form II - Canonic:}
Direct Form II shares the delay elements between the feedforward and feedback paths. We introduce an intermediate variable $w[n]$.
$$w[n] = x[n] + 1.1w[n-1] - 0.3w[n-2]$$
$$y[n] = w[n] + 0.5w[n-1]$$

\textbf{(b) Memory Comparison:}
\begin{itemize}
    \item \textbf{Direct Form I:} Requires $M + N = 1 + 2 = 3$ delay elements ($x[n-1]$, $y[n-1]$, $y[n-2]$).
    \item \textbf{Direct Form II:} Requires $\max(M, N) = \max(1, 2) = 2$ delay elements ($w[n-1]$, $w[n-2]$).
    \item \textbf{Conclusion:} Direct Form II is more memory efficient (Canonical).
\end{itemize}

\textbf{(c) Numerical Differences:}
Direct Form II is more sensitive to quantization effects (overflow/round-off) at the internal nodes ($w[n]$) compared to Direct Form I, especially for high-Q poles, because the poles solely determine the magnitude of the intermediate variable $w[n]$ before zeros can attenuate it.

\begin{figure}[h]
    \centering
    \includegraphics[width=0.8\textwidth]{filter_comparison.png}
    \caption{Filter Output Comparison}
    \label{fig:filter_comp}
\end{figure}

\appendix
\section{Python Source Code}
\textbf{File:} \texttt{filter\_realization.py}
\begin{lstlisting}[caption={Filter Realization: Direct Form I vs Scipy lfilter -- Assignment 6}]
import numpy as np
import matplotlib.pyplot as plt
from scipy.signal import lfilter

def manual_filter(b, a, x):
    """
    Implement LCCDE manually:
    a[0]y[n] = b[0]x[n] + b[1]x[n-1] + ... - a[1]y[n-1] - ...
    Assumes a[0] = 1.
    """
    y = np.zeros_like(x)
    M = len(b)
    N = len(a)
    for n in range(len(x)):
        for i in range(M):
            if n - i >= 0:
                y[n] += b[i] * x[n - i]
        for j in range(1, N):
            if n - j >= 0:
                y[n] -= a[j] * y[n - j]
    return y

def main():
    b = [1, 0.5]
    a = [1, -1.1, 0.3]
    fs = 100
    t = np.arange(0, 1, 1/fs)
    x = np.sin(2*np.pi*5*t) + 0.5*np.sin(2*np.pi*20*t) + 0.1*np.random.randn(len(t))
    
    y_manual = manual_filter(b, a, x)
    y_lfilter = lfilter(b, a, x)
    difference = y_manual - y_lfilter
    max_diff = np.max(np.abs(difference))
    print(f"Max difference between Manual and Lfilter: {max_diff:.5e}")
    
    plt.figure(figsize=(10, 6))
    plt.subplot(2, 1, 1)
    plt.plot(t, x, label='Input x[n]', alpha=0.7)
    plt.plot(t, y_manual, label='Output y[n] (Manual)', linestyle='--')
    plt.title('Filter Implementation Verification')
    plt.legend()
    plt.grid(True)
    plt.ylabel('Amplitude')
    
    plt.subplot(2, 1, 2)
    plt.plot(t, difference, label='Difference (Manual - Lfilter)', color='r')
    plt.title(f'Numerical Difference (Max Error: {max_diff:.2e})')
    plt.xlabel('Time (s)')
    plt.ylabel('Error')
    plt.grid(True)
    plt.legend()
    plt.tight_layout()
    plt.savefig('filter_comparison.png')
    print("Plot saved to filter_comparison.png")

if __name__ == "__main__":
    main()
\end{lstlisting}

\end{document}
